\documentclass[11pt]{article}
\usepackage[margin=1in]{geometry}
\usepackage{amsmath,amssymb,amsthm,hyperref}
\newtheorem{theorem}{Theorem}
\newcommand{\Pf}{\operatorname{Pf}}
\title{Equal Pfaffian invariants in distinct symplectic congruence orbits}
\author{Formal verification of a posted counterexample by Shivam Patel}
\date{October 26, 2026}
\begin{document}\maketitle
\begin{abstract}
We verify a six-dimensional counterexample to the symplectic-form orbit-classification conjecture recorded as MathDB 381863. Two nonsingular skew-symmetric forms, neither a scalar multiple of the standard form, have every proposed invariant equal but lie in different symplectic congruence orbits. The finite Pfaffian formulas and the complete orbit obstruction are proved in Lean. The matrices are from Shivam Patel's existing posted counterexample; no discovery or priority claim is made.
\end{abstract}
The actual primary source defines
\[
 S(2n,\mathbb R)=\{A:A^T=-A,\ \det A\ne0\},\qquad
 J=\operatorname{diag}(J_2,\ldots,J_2),\quad
 J_2=\begin{pmatrix}0&1\\-1&0\end{pmatrix},
\]
and $\operatorname{Sp}(2n)=\{P:P^TJP=J\}$, acting by $A\mapsto P^TAP$. For $0\le k<n$, its invariant $s_k(A)$ is the sum of the Pfaffians of the principal submatrices obtained by deleting $k$ complete coordinate pairs. Section 9, Conjecture 1 asserts that these invariants determine the orbit for matrices that are not scalar multiples of $J$.

The word ``symmetric'' in MathDB's description of $S(2n,\mathbb R)$ is a transcription error. The present theorem concerns the actual primary-source symplectic forms, which are skew-symmetric.
\begin{theorem}
Let $n=3$, let $J=\operatorname{diag}(J_2,J_2,J_2)$, and set
\[
A_0=\begin{pmatrix}
0&1&0&0&0&0\\-1&0&0&0&0&0\\0&0&0&1&0&0\\
0&0&-1&0&0&0\\0&0&0&0&0&2\\0&0&0&0&-2&0
\end{pmatrix},\quad
A_1=\begin{pmatrix}
0&1&0&0&0&0\\-1&0&-1&0&0&0\\0&1&0&1&0&0\\
0&0&-1&0&0&0\\0&0&0&0&0&2\\0&0&0&0&-2&0
\end{pmatrix}.
\]
Both matrices lie in $S(6,\mathbb R)$ and neither is scalar $J$. They satisfy $s_k(A_0)=s_k(A_1)$ for $k=0,1,2$, but there is no symplectic matrix $P$ with $A_1=P^TA_0P$.
\end{theorem}
\begin{proof}
Skewness follows directly from the entries. The two Pfaffians of order six are both $2$, so both determinants are $4$; independently, the formal certificate proves nonvanishing by explicit matrix inverses. The entries in positions $(1,2)$ and $(5,6)$ exclude $A_0=cJ$, and the entry $(2,3)=-1$ excludes $A_1=cJ$.

To see equality of all the Pfaffian invariants directly, note that the only additional edge in the upper-triangular part of $A_1$ is between coordinates 2 and 3. In any matching contribution using that edge, coordinate 1 would need a partner other than 2, but all its corresponding entries are zero. Hence that contribution vanishes. This proves equality of the full Pfaffians. For the four-coordinate submatrices obtained by deleting one complete pair, deletion of either of the first two pairs removes the added edge, while deletion of the last pair leaves the same matching obstruction. Finally the two-coordinate pair submatrices are identical. Thus all three invariant sums coincide.

Put $K_i=-JA_i$. Since $J^2=-I$, congruence by a symplectic $P$ would imply similarity of the $K_i$. More explicitly, set $L=-JP^TJ$. The symplectic identity gives $LP=PL=I$, and the congruence identity gives $K_1=LK_0P$. Direct arithmetic gives
\[
 (K_0-I)(K_0-2I)=0,\qquad (K_1-I)(K_1-2I)\ne0.
\]
Indeed $K_0=\operatorname{diag}(1,1,1,1,2,2)$, whereas the $(1,3)$ entry of the latter product is $-1$. But similarity would yield
\[
 (K_1-I)(K_1-2I)
 =L(K_0-I)(PL)(K_0-2I)P=0,
\]
a contradiction. This proves the orbit obstruction.
\end{proof}
For clarity about normalization, the ordinary Pfaffian formula on a $2m$-by-$2m$ skew matrix is
\[
 \Pf(M)=\frac1{2^m m!}\sum_{\sigma\in S_{2m}}\operatorname{sgn}(\sigma)
       \prod_{i=1}^m M_{\sigma(2i-1),\sigma(2i)}.
\]
Printed Definition 2.12 omits the prefactor while saying the sum ranges over all permutations. Our certificate verifies invariant equality for both the raw printed sum and the standard normalized formula. For each $k$, every relevant submatrix has the same size, so the two invariant conventions differ only by the common factor $2^{3-k}(3-k)!$; equality is unaffected.

The module \texttt{SymplecticOrbits.lean} defines this full permutation sum and explicitly enumerates every complete-pair deletion for $n=3$. The lemmas \texttt{invariant\_zero}, \texttt{invariant\_one}, and \texttt{invariant\_two} prove the exact bridge to the three source sums of normalized Pfaffians. All three integer equalities are checked by Lean's kernel, and proved cast identities transport the finite sums to the displayed real matrices. This uses actual source invariants, with no assumed list of spectral or polynomial coefficients. The remaining declarations prove both nonsingularity conditions, both non-scalar conditions, and the complete congruence-to-similarity obstruction. The final theorem is \texttt{SymplecticOrbitCounterexample.source\_counterexample}.

This is verification of the matrices in Patel's posted counterexample. The local proof replaces an eigenspace-dimension computation by an explicit quadratic-polynomial obstruction. No new-to-literature or first-formalization claim is made. The full primary source, post identifiers, and content hashes accompany the certificate.
\begin{thebibliography}{9}
\bibitem{source} L. Shi, S. Joshi, and R. Bhargava, \emph{Congruence of Linear Symplectic Forms by the Symplectic Group}, arXiv:2212.08360v1, Section 1.2, Definition 2.12, Theorem 2.8, and Section 9, Conjecture 1. \url{https://arxiv.org/abs/2212.08360v1}.
\bibitem{post} S. Patel, posted counterexample to MathDB 381863. \url{https://mathdb.com/p/381863}.
\end{thebibliography}
\end{document}
